\section{System Model}
\label{sec:ch3-system-model}

\subsection{Conceptual Framework}
\label{sec:ch3-conceptual-framework}

The prototype was modeled as an input--process--output system. Its inputs were
the mains-voltage sensor, current sensor, thermistor, and wireless connection;
its processing stage was the ESP32-S3 firmware with calibration, thermal
estimation, waveform-feature extraction, TinyML inference, relay control, and
cloud logging; and its outputs were warnings, PWA monitoring, stored event
records, and relay disconnection when a protection condition was latched.

Figure~\ref{fig:ch3-conceptual-framework} summarizes this research model.
\begin{figure}[H]
  \centering
  \begin{tikzpicture}[node distance=0.32in and 0.30in]
    \node[thesis process,text width=1.50in] (inputs) {Inputs\\voltage, current, device temperature, network state};
    \node[thesis process,text width=1.65in,right=of inputs] (processing) {Embedded processing\\calibration, thermal model, waveform features, TinyML inference};
    \node[thesis process,text width=1.50in,right=of processing] (outputs) {Outputs\\warnings, relay action, logs, PWA display/control};
    \node[thesis process,text width=1.45in,below=0.34in of processing] (validation) {Controlled validation\\sensor, overload, socket heating, series arc};
    \node[thesis process,text width=1.45in,below=0.34in of validation] (interpretation) {Interpretation\\complementary protection under tested conditions};
    \draw[thesis arrow] (inputs) -- (processing);
    \draw[thesis arrow] (processing) -- (outputs);
    \draw[thesis arrow] (validation) -- (processing);
    \draw[thesis arrow] (processing) -- (interpretation);
    \draw[thesis arrow] (outputs.south) |- (interpretation.east);
  \end{tikzpicture}
  \caption{Conceptual framework of the TinyML Smart Plug.}
  \label{fig:ch3-conceptual-framework}
\end{figure}

Figure~\ref{fig:ch3-conceptual-framework} shows that the study did not treat
the smart plug as a measurement-only device. The sensing inputs are converted
into protection variables and waveform features, then used to support user
notification, fault logging, and local relay action. This distinction is
important because the prototype was intended as complementary protection, not
as a replacement for branch-circuit breakers, AFCIs, GFCIs, wiring inspection,
or certified protective equipment.

\subsection{System Architecture}
\label{sec:ch3-system-architecture-discussion}

The same behavior can be described as a layered embedded system. The sensing
layer captures voltage, current, and socket temperature; the firmware layer
computes calibrated variables and protection states; the TinyML layer provides
context-aware series-arc inference; and the PWA layer displays live and
historical data using Firebase services [55], [56].

Figure~\ref{fig:ch3-system-architecture} presents the system architecture used
to connect these layers.
\begin{figure}[H]
  \centering
  \begin{tikzpicture}[node distance=0.24in and 0.27in]
    \node[thesis process,text width=1.20in] (sensors) {Sensors\\voltage, current, temperature};
    \node[thesis process,text width=1.55in,right=of sensors] (esp) {TinyML Smart Plug\\ESP32-S3 local protection};
    \node[thesis process,text width=1.28in,right=of esp] (relay) {Relay and\\protected outlet};
    \node[thesis process,text width=1.28in,above=0.30in of relay] (local) {OLED, buzzer,\\local indicators};
    \node[thesis process,text width=1.38in,below=0.32in of esp] (firebase) {Firebase\\Realtime Database};
    \node[thesis process,text width=1.42in,right=of firebase] (pwa) {PWA\\visualization, logs, basic control};
    \draw[thesis arrow] (sensors) -- (esp);
    \draw[thesis arrow] (esp) -- node[above,font=\scriptsize]{local trip} (relay);
    \draw[thesis arrow] (esp) -- (local);
    \draw[thesis arrow] (esp) -- node[left,font=\scriptsize]{telemetry} (firebase);
    \draw[thesis arrow] (firebase) -- node[above,font=\scriptsize]{live data} (pwa);
    \draw[thesis arrow] (pwa.south west) -- node[below,font=\scriptsize]{clear/relay commands} (firebase.south east);
    \draw[thesis arrow] (firebase.north) -- node[left,font=\scriptsize]{control fields} (esp.south);
  \end{tikzpicture}
  \caption{System architecture of the TinyML Smart Plug.}
  \label{fig:ch3-system-architecture}
\end{figure}

Figure~\ref{fig:ch3-system-architecture} places the ESP32-S3 at the center of
the design. The microcontroller receives conditioned analog signals, runs the
protection and TinyML routines locally, drives the relay and indicators, and
uploads selected telemetry. Firebase and the PWA also carry basic remote
control fields, such as fault clearing and relay-control commands, but these
commands do not replace the local protection logic. Keeping the protection
decision on the
microcontroller prevents the relay action from depending on internet
availability, while the cloud and PWA remain useful for visualization,
logging, labeling, review, and operator interaction.

\subsection{Input, Processing, and Output Summary}
\label{sec:ch3-system-ipo}

The design variables used by the system are summarized in
Table~\ref{tab:ch3-system-variables}. The table separates direct measurements
from computed variables so that later equations can be traced to the firmware
implementation.
\begin{table}[H]
  \centering
  \caption{Main firmware variables used by the prototype}
  \label{tab:ch3-system-variables}
  \begin{tabular}{p{0.25\textwidth}p{0.28\textwidth}p{0.35\textwidth}}
    \toprule
    \textbf{Variable group} & \textbf{Source} & \textbf{Application} \\
    \midrule
    RMS voltage & ZMPT101B channel and voltage calibration & Voltage monitoring, mains-present logic, undervoltage and overvoltage protection \\
    RMS current & ACS37035 channel, MCP3204 ADC, anti-aliasing filter, and current calibration & Load monitoring, overload protection, thermal load term, and arc-feature extraction \\
    Device temperature measurement & 10k$\Omega$ thermistor and divider computation & Baseline thermal measurement for socket-temperature estimation \\
    Waveform features & Current samples processed by the arc-detection module & Context classification and Random Forest arc inference \\
    Protection state & Protection manager and relay driver & Warning, latching, relay disconnection, and event logging \\
    \bottomrule
  \end{tabular}
\end{table}

Table~\ref{tab:ch3-system-variables} shows that current is used by several
subsystems: it supports ordinary overload detection, thermal compensation, and
the TinyML feature vector. Therefore, the current channel was treated as a
shared measurement path whose calibration and sampling behavior must remain
consistent across protection functions.
